% Title block, two-sentence abstract, and scope note. The availability and
% deflation housekeeping is a footnote on the author line rather than a side
% column, so the abstract runs the full text block.
% The scope note is load-bearing: it tells a reader what question the study
% answers and, just as importantly, which arguments it declines to enter.
\thispagestyle{empty}

\begin{bannerbox}
\color{white}
{\fontsize{15}{18}\selectfont\sffamily\bfseries Economic Opportunity for Musicians and Creative Workers:\\Austin in State and National Context}\\[4pt]
{\fontsize{11.5}{14}\selectfont\sffamily Earnings, venues, public support, 2001 to 2026}\\[5pt]
{\footnotesize\color{amlightbg}Working Paper \textbar\ August 2026}
\end{bannerbox}

\vspace{7pt}

\noindent{\small\textbf{Eric A. Booth}}\footnote{The analysis code, source
data, and full evidence file behind every figure in this study are available
from the author. All dollar figures are inflation-adjusted to 2025 dollars
unless noted.}\\[1pt]
{\footnotesize\color{ammuted}\href{mailto:eric.a.booth@gmail.com}{eric.a.booth@gmail.com}}

\vspace{8pt}

\noindent This study gathers the public evidence on the economic position of
musicians and creative workers in Austin and compares it with the city, state
and federal programs directed at them. Austin musicians report median earnings
from work of \NPumsAustinMusMedpernpPos\ against
\NPumsAustinAllworkMedpernpPos\ for all Austin workers, about half, and
\NPumsAustinMusSelfemp\ of them work for themselves, which places much of this
workforce at or beyond the edge of the payroll statistics that public programs
are designed around (see Section~\ref{sec:payroll}). Musicians occupy the
low-earnings end of a wider creative workforce this study
also examines. Across nine pooled creative occupations,\footnote{Musicians and
singers; music directors and composers; graphic designers; photographers;
media, audio-visual and sound-equipment workers; producers and directors;
actors; dancers and choreographers; writers and authors. These are the American
Community Survey occupation codes that isolate creative and cultural workers in
the Texas public use microdata; sound engineers are inside the combined media
category and cannot be separated. This study reports the pooled figure and the
musician figure separately, never one inside the other.} the Austin median
rises to \NPumsAustinCreativeMedpernpPos\ and
\NPumsAustinCreativeSelfemp\ work for themselves. The report is mainly about
musicians because that is where the measurement record runs deepest, the
earnings gap runs widest and the public programs are most concentrated; it
turns to the wider creative workforce for occupation comparisons in
Section~\ref{sec:earnings} and for the state's arts-economy aggregate in
Section~\ref{sec:state}.

\vspace{9pt}

\begin{keyfinding}
{\small\textbf{What this study does.} It assembles what public data can measure
about the economic position of people who work in music in Austin, and sets that
record against the city and state programs aimed at them. It draws on
federal statistics, city and state administrative records, and an original panel
of monthly alcohol-sales receipts for \NVenuePanelIsCurated, a curated set of
named Austin live-music rooms rather than a census of them, covering 2007
through mid-2026 (see Section~\ref{sec:venues}).

\smallskip
\textbf{What it does not do.} It does not take a position in the running
argument over whether musicians are paid fairly, or over what a venue owes a
performer. That argument turns on judgments about value and risk that no dataset
can settle, and the measurement record is too thin to resolve it either way.
This study stays with what is traceable: earnings recorded in federal surveys,
receipts recorded on tax filings, and dollars recorded in public budgets.}
\end{keyfinding}

\vspace{4pt}
